\subsection{Morphisms of principal fibrations}

6.3.1. Let \(\lambda = (\mathbf{P}, \mathbf{G}, \mathbf{B}, \pi)\) and \(\lambda' = (\mathbf{P}', \mathbf{G}', \mathbf{B}', \pi')\) be two principal fibrations. A morphism of \(\lambda\) to \(\lambda'\) is a triple \((f, \varphi, h)\) , where \(f: \mathbf{P} \to \mathbf{P}'\) and \(h: \mathbf{B} \to \mathbf{B}'\) are morphisms, \(\varphi: \mathbf{G} \to \mathbf{G}'\) is a homomorphism of group varieties, and where \(\pi' \circ f = h \circ \pi\) and \(f(x \cdot g) = f(x) \cdot \varphi(g)\)

for \(x \in \mathbf{P}, g \in \mathbf{G}\) . Note that \(f\) determines \(h\) ; one will often say that \((f, \varphi)\) , or even simply \(f\) is a morphism.

When \(B = B'\) and \(h = \mathrm{Id}_{B}\) (resp. when \(G = G'\) and \(\varphi = \mathrm{Id}_{G}\) ), a morphism is called a B-morphism compatible with \(\varphi\) (resp. a G-morphism compatible with h). A morphism that is both a B-morphism and a G-morphism is called a G-B-morphism (here again, one often says "morphism" when there can be no confusion). Every G-B-morphism \(f : P \to P'\) is an isomorphism of the variety P onto the variety \(P'\) ; the inverse isomorphism \(f^{-1}\) is a G-B-morphism: f is a G-B-isomorphism of P onto \(P'\) .

Two principal fiber bundles P and P' with the same base B and the same structural group G are said to be G-B-isomorphic (or simply isomorphic) if there exists a G-B-isomorphism of P onto P'.

6.3.2. The principal fibration \((\mathbf{B} \times \mathbf{G}, \mathbf{G}, \mathbf{B}, \mathbf{pr}_{1})\) is called the trivial principal fibration with base B and structural group G. An isomorphism of a principal fibration \(\lambda = (\mathbf{P}, \mathbf{G}, \mathbf{B}, \pi)\) onto the trivial principal fibration with base B and structural group G is called a trivialization of \(\lambda\) . Every section s of \(\lambda\) defines a trivialization \(f_{s}\) of \(\lambda\) by the formula:

\[
f _ {s} ^ {- 1} (b, g) = s (b). g \quad \text {   for   } b \in \mathbf {B} \text {   and   } g \in \mathbf {G}.
\]

Thus we obtain a bijection from the set of sections of \(\lambda\) onto the set of trivializations of \(\lambda\) . Moreover, if \(s_0\) is a section of \(\lambda\) , every section \(s\) of \(\lambda\) can be written uniquely in the form \(s(b) = s_0(b) \cdot r(b)\) , where \(r: B \to G\) is a morphism.

6.3.3. Let \(\lambda = (\mathbf{P}, \mathbf{G}, \mathbf{B}, \pi)\) be a principal fibration, and let \(h: \mathbf{B}' \to \mathbf{B}\) be a morphism. Let \(\pi'\) (resp. \(h'\) ) be the canonical morphism from \(\mathbf{B}' \times_{\mathbf{B}} \mathbf{P}\) to \(\mathbf{B}'\) (resp. into \(\mathbf{P}\) ). Let us make \(\mathbf{G}\) act on \(\mathbf{B}' \times_{\mathbf{B}} \mathbf{P}\) by the formula

\[
(b ^ {\prime}, x). g = (b ^ {\prime}, x. g), \quad (b ^ {\prime}, x) \in \mathbf {B} ^ {\prime} \times_ {\mathbf {B}} \mathbf {P}, \quad g \in \mathbf {G}.
\]

The quadruple \((\mathbf{B}' \times_{\mathbf{B}} \mathbf{P}, \mathbf{G}, \mathbf{B}', \pi')\) is a principal fibration, called the inverse image of \(\lambda\) by \(h\) , and denoted \(\mathbf{B}' \times_{\mathbf{B}} \lambda\) or also \(h^*\lambda\) . The map \(h': \mathbf{B}' \times_{\mathbf{B}} \mathbf{P} \to \mathbf{P}\) is a G-morphism compatible with \(h\) ; it enjoys the following universal property: if \(f\) is a G-morphism compatible with \(h\) of a principal fibration \(\lambda'\) with base \(\mathbf{B}'\) into the fibration \(\lambda\) , there exists a unique G-B'-isomorphism \(k: \lambda' \to \mathbf{B}' \times_{\mathbf{B}} \lambda\) such that \(f = h' \circ k\) .

When B' is a subvariety of B, and h is the canonical injection of B' into B, \(B' \times_{B} P\) is identified with the submanifold \(\pi^{-1}(B')\) of P; it is called the principal fiber bundle induced by P over B', and it is denoted \(\pi^{-1}(B')\) , or equivalently P\textbar B'. Every \(x \in B\) has an open neighborhood U such that P\textbar U is trivial.

